\section{Historischer Abriss und Originalarbeiten}

Die Darstellung verbindet zwei ältere Fragen: Wie organisiert man
wiederholte Additionen, und wie gewinnt man aus einer Wertetafel
Zwischenwerte? Das erste Problem führt zum arithmetischen Dreieck,
das zweite zur Interpolation. Die heutige Formel fasst diese Gedanken
in einer einheitlichen Schreibweise zusammen.

\textbf{Pascal: die Arithmetik des Dreiecks.}
Pascals 1665 veröffentlichter \emph{Traité du triangle arithmétique}
entwickelt die Bildung der Einträge durch Addition zweier Vorgänger
und ihre Summationsgesetze. Der Definitionsabschnitt sowie die
Conséquences II--III zeigen damit die Rechenstruktur hinter (P) und
(A) \cite{newton:pascal-source}. Das ist eine wichtige systematische
Darstellung, keine Behauptung, das Zahlendreieck sei erst bei Pascal
entstanden. Unsere Argumentation benötigt seine Additionsregel in
moderner Binomialschreibweise.

\textbf{Newton: Interpolation aus gegebenen Werten.}
Im fünften Lemma des dritten Buches der \emph{Principia} von 1687
beschreibt Newton die Konstruktion einer Kurve durch vorgegebene
Punkte: Der erste Fall benutzt gleiche Abstände, der zweite ungleiche
Abstände und entsprechend angepasste Differenzen. Schon das folgende
Lemma wendet das Verfahren auf die Bestimmung von Kometenpositionen
aus Beobachtungsdaten an \cite[Buch~III, Lemmata~V--VI]{newton:principia}.
Die Differenzentafel dient hier der Berechnung unbekannter Werte aus
bekannten Werten.

Newtons 1711 veröffentlichte \emph{Methodus differentialis} behandelt
die Interpolation ausführlicher, auch mit zentral angeordneten
Stützstellen. Die Propositionen I--III entwickeln Differenzen und
Interpolationsformen; die späteren Propositionen und das Scholium
verbinden sie mit Tabellenberechnung und Quadratur
\cite{newton:methodus}. Das historische Problem ist somit breiter
als unser Satz über Folgen. Die spezielle Darstellung an
$0,1,2,\ldots$ ist eine besonders einfache Form dieses Verfahrens.

\textbf{Taylor: Differenzenrechnung und der Übergang zu Ableitungen.}
In der \emph{Methodus incrementorum directa \& inversa} von 1715
leitet Taylor in Proposition VII eine Entwicklung aus der
Differenzentafel ab. Das zweite Korollar verbindet sie mit der
entsprechenden Entwicklung in Fluxionen, also in der damaligen Sprache
mit Ableitungen. Das Scholium zu Proposition VIII nennt den Abbruch
bei polynomialen Ausdrücken ausdrücklich
\cite[Pars~I, S.~21--24]{newton:taylor}.
Damit begegnen dort beide Leitgedanken dieser Lesefassung: die
Entwicklung aus Anfangsdifferenzen und ihr Ende bei Polynomen.

Die Bezeichnungen $\Delta$, $B_k$ und die Operatorenschreibweise sind
hier moderne Formulierungen. Die genannten Stellen dokumentieren
konkrete Beiträge und Anwendungen; sie begründen keine ausschließliche
Priorität Newtons für die gesamte Interpolationstheorie. Auch ist die
vorliegende Induktion eine heutige knappe Herleitung und keine
Übersetzung eines historischen Beweises.

\begin{thebibliography}{9}
\small
\setlength{\itemsep}{.65em}

\bibitem{newton:pascal-source}
Blaise Pascal: \emph{Traité du triangle arithmétique, avec quelques
autres petits traitez sur la mesme matière}.
Paris: Guillaume Desprez, 1665.
Hier: Definitionsabschnitt, S.~1--2, und Conséquences II--III.
\href{https://gallica.bnf.fr/ark:/12148/btv1b86262012}{Originaldigitalisat der Bibliothèque nationale de France}.
Eine \href{https://fr.wikisource.org/wiki/Trait%C3%A9_du_triangle_arithm%C3%A9tique}{Transkription des Textes}
ist ebenfalls zugänglich; sie ersetzt nicht die historische Druckausgabe.

\bibitem{newton:principia}
Isaac Newton: \emph{Philosophiae naturalis principia mathematica}.
London, 1687. Buch~III, Lemma~V, Druckseiten~481--483;
Lemma~VI, Druckseite~483.
\href{https://www.newtonproject.ox.ac.uk/view/texts/normalized/NATP00301}{Lateinischer Originaltext in der wissenschaftlichen Transkription des Newton Project}.
Die Seitenangaben beziehen sich auf die Ausgabe von 1687.

\bibitem{newton:methodus}
Isaac Newton: \emph{Methodus differentialis}.
In: William Jones (Hg.), \emph{Analysis per quantitatum series,
fluxiones, ac differentias: cum enumeratione linearum tertii ordinis}.
London, 1711, S.~93--101. Hier insbesondere Propositionen~I--III,
S.~93--97, sowie Propositionen~V--VI und Scholium, S.~99--101.
\href{https://www.e-rara.ch/download/pdf/2621379.pdf}{Originaldigitalisat des Abschnitts};
\href{https://doi.org/10.3931/e-rara-8934}{dauerhafter Nachweis des Gesamtwerks bei e-rara / ETH-Bibliothek}.

\bibitem{newton:taylor}
Brook Taylor: \emph{Methodus incrementorum directa \& inversa}.
London, 1715. Pars~I, Propositio~VII / Theorema~III, S.~21--23;
Corollarium~II, S.~23; Propositio~VIII, Scholium, S.~24.
\href{https://archive.org/details/UFIE003454_TO0324_PNI-2529_000000}{Originaldigitalisat der Ausgabe von 1715}.
Zum leichteren Lesen:
\href{https://www.17centurymaths.com/contents/brooktaylor/methodofincrementsparttwo.pdf}{Ian Bruces Ausgabe mit lateinischem Text und englischer Übersetzung, Teil I~B}.
Die Fundstellen beziehen sich auf die Druckseiten des Originals.

\end{thebibliography}
